\begin{scanpage}{0057}{40}
\[\left.\begin{aligned}E_z=-H_\phi&=E_{0z}\left\{-\frac{\cos[ka(\sin\alpha-\sin\phi)]}{\cos\frac{\alpha+\phi}{2}}+\right.\\&\qquad\left.+i\frac{\sin[ka(\sin\alpha-\sin\phi)]}{\sin\frac{\alpha-\phi}{2}}\right\}\frac{e^{i(kr+\frac\pi4)}}{\sqrt{2\pi kr}},\\E_\phi=H_z&=H_{0z}\left\{\frac{\cos[ka(\sin\alpha-\sin\phi)]}{\cos\frac{\alpha+\phi}{2}}+\right.\\&\qquad\left.+i\frac{\sin[ka(\sin\alpha-\sin\phi)]}{\sin\frac{\alpha-\phi}{2}}\right\}\frac{e^{i(kr+\frac\pi4)}}{\sqrt{2\pi kr}}.\end{aligned}\right\}\tag{6.14}\]
These equations are valid when $r\gg ka^2$ and $|\phi|<\frac\pi2$. Moreover, it is assumed that $ka\gg1$, since only under this condition is one able to consider the nonuniform part of the current in the vicinity of the strip's edge to be approximately the same as on the corresponding half-plane. In the case of normal incidence of a plane wave $(\alpha=0)$, Equations (6.14) change into expressions corresponding to the first approximation of Schwarzschild [15].

From relationships (6.03) and (6.05), it follows that the electric field is an even function, and the magnetic field an odd function, of the $x$ coordinate measured perpendicular to the plane $x=0$ (in which the current flows)
\[E_z(x)=E_z(-x),\quad H_z(x)=-H_z(-x).\tag{6.15}\]
Therefore, on the basis of Equations (6.14) and (6.15) one is able to write the expressions for the fringing field in the region $x<0$ (where $\frac\pi2<|\phi|<\pi$)
\[\left.\begin{aligned}E_z=-H_\phi&=\pm E_{0z}\left\{\frac{\cos[ka(\sin\alpha-\sin\phi)]}{\sin\frac{\alpha-\phi}{2}}-\right.\\&\qquad\left.-i\frac{\sin[ka(\sin\alpha-\sin\phi)]}{\cos\frac{\alpha+\phi}{2}}\right\}\frac{e^{i(kr+\frac\pi4)}}{\sqrt{2\pi kr}},\\E_\phi=H_z&=\pm H_{0z}\left\{\frac{\cos[ka(\sin\alpha-\sin\phi)]}{\sin\frac{\alpha-\phi}{2}}+\right.\end{aligned}\right\}\tag{6.16}\]
\end{scanpage}
